\subsection{Tangens: Periodisch, aber nicht beschränkt}
\label{subsec:real_functions_tangent}
Der Tangens ist der Quotient aus Sinus und Kosinus. Seine
Definitionsmenge folgt deshalb aus den Nullstellen des Kosinus.
\begin{DefinitionBox}{Tangensfunktion}{real_functions_tangent}
  Die Tangensfunktion ist definiert durch
  \[
    \tan x=\frac{\sin x}{\cos x},
    \qquad D=\R\setminus\left\{\frac\pi2+k\pi\mid k\in\Z\right\}.
  \]
  Sie ist ungerade, auf $D$ stetig und hat die Grundperiode $\pi$.
\end{DefinitionBox}
\begin{center}
  \CachedTikzFigure{content/resources/chapters/real_functions/tikz/functions_tangent.tex}
\end{center}
\begin{KeyPointBox}{Jeder Tangensast durchläuft alle reellen Werte}{real_functions_tangent_branch}
  Auf jedem Intervall $(-\pi/2+k\pi,\pi/2+k\pi)$ ist der Tangens
  streng monoton steigend und nimmt jeden reellen Wert genau einmal an.
  Für den mittleren Ast gilt
  \[
    \lim_{x\to(-\pi/2)^+}\tan x=-\infty,\qquad
    \lim_{x\to(\pi/2)^-}\tan x=+\infty.
  \]
  Die Stellen $\pi/2+k\pi$ sind Polstellen, nicht Stellen mit einem
  Funktionswert \glqq unendlich\grqq.
\end{KeyPointBox}

\clearpage
\subsection{Harmonische Schwingungen}
\label{subsec:real_functions_harmonic_oscillations}
Mit den bekannten Graphentransformationen passen wir den Sinus an
eine Schwingung mit vorgegebener Auslenkung und Dauer an.
\par\noindent\begin{minipage}{\linewidth}
\begin{DefinitionBox}{Harmonische Schwingung}{real_functions_harmonic}
  Eine \emph{harmonische Schwingung} \EnglishTerm{harmonic oscillation}
  lässt sich mit $A>0$ und $\omega>0$ schreiben als
  \[
    y(t)=y_0+A\sin\bigl(\omega(t-t_0)\bigr).
  \]
  Hier bedeuten $y_0$ die Mittellage, $A$ die \emph{Amplitude}
  \EnglishTerm{amplitude}, $\omega$ die \emph{Kreisfrequenz}
  \EnglishTerm{angular frequency} und $t_0$ die zeitliche Verschiebung.
  Die \emph{Periodendauer} \EnglishTerm{period duration} und die
  \emph{Frequenz} \EnglishTerm{frequency} sind
  \[
    T=\frac{2\pi}{\omega},\qquad \nu=\frac1T.
  \]
\end{DefinitionBox}
\begin{ExampleBox}[before skip=0.3\baselineskip]{Eine Schwingung aus dem Diagramm beschreiben}{real_functions_oscillation_example}
  Die dargestellte Auslenkung wird in Millimetern, die Zeit in Sekunden
  gemessen. Bestimmen Sie $A$, $T$, $\omega$ und einen passenden Term.
  \begin{center}
    \CachedTikzFigure{content/resources/chapters/real_functions/tikz/functions_oscillation.tex}
  \end{center}
  \begin{ClozeArea}[34]{6}
    Die Mittellage ist null; bei $t=0$ beginnt die Bewegung mit
    steigender Auslenkung. Damit gilt
    \[
      A=2\,\mathrm{mm},\quad T=4\,\mathrm{s},\quad
      \omega=\frac{\pi}{2}\,\mathrm{s}^{-1},\qquad
      y(t)=2\,\mathrm{mm}\,\sin\left(\frac{\pi t}{2\,\mathrm{s}}\right).
    \]
  \end{ClozeArea}
\end{ExampleBox}
\end{minipage}

In der Schreibweise $A\sin(\omega t+\varphi)$ heißt
$\varphi=-\omega t_0$ die \emph{Phase}
\EnglishTerm{phase}; auch dieses Winkelargument ist dimensionslos.
